\section{Persistence, the trigger, pool occupancy and goodput}\label{app:dataonly}

Four analyses (D1, D2, D5 and the pool occupancy) use only runs committed before 2026-11-02; a fifth, D4, a post-hoc
reading of Lu et al.'s model on our cells, is in Appendix~\ref{app:lu}. Each was computed by one agent and then recomputed by a second agent
with its own code from the raw files; where they differ we report the checked value and say so. They are post hoc. Scripts and
outputs: \fpath{analysis/r15_dataonly/} and \fpath{results/r15_dataonly/} (\fpath{PROVENANCE.md} there).

\paragraph{Persistence (D5).} Rolling 300\,s windows of the follow-up hit fraction; a trap starts at the first window after the
trigger with hit fraction at most 0.2, and an escape is any later window with at least 0.5 before the run ends or an intervention
begins (a rescue cap, a SHED stop). In the 50 trapped FCFS runs committed before the blind tests of 2026-11-05 (49 after a burst, registered and exploratory,
and the trapped run of the 900\,s TOOLSLOW arm; the TOOLSLOW-W and CUT runs are not included) there was no escape; the highest window after onset reached 0.164. Trapped time totals 99--103\,h, depending on whether the
trap is taken to start at the beginning or the end of its first window (the first computation gave 104.5\,h, which the recomputation
did not reproduce), and the shortest trapped span is 6100--6400\,s, so a one-sided 95\% upper bound on a run's escape probability
within such a span is $1-0.05^{1/50}=5.8$\%. Stock LPM escaped in 4 of its 9 trapped runs (several barely above 0.5), and LPM with its
fallback disabled in its one trapped run. Of the 16 from-empty runs that collapsed, one recovered within its 3000\,s run (150k pool,
1.30~$\lref$, seed~0).

\paragraph{The trap outlives its trigger (D1).} In 16 stuck runs committed before those tests, every trigger session (the burst's, or for TOOLSLOW the sessions
with a slowed tool call) had left the system 558--6246\,s before the end: the 600k H200 paired test (6 of 6), BASIN bursts of 50 (1),
60 (3) and 100 (3), SHED's 600\,s stop (2) and TOOLSLOW's trapped run (1). The 727--15{,}453 turns completed after the last trigger
session left had no hit in 15 of these runs; in the TOOLSLOW run 4 of 5{,}937 hit, all within 715\,s after its last slowed session left. In
the 24 registered 200-session runs on the A100 (SGLang 300k and native, vLLM, Llama) 3--30 burst sessions were still present at the
end (at most 4.1\% of the sessions present; the first computation's ``5.8\%'' could not be reproduced), so for those runs we do not
claim that the trap outlived its trigger.

\paragraph{Pool occupancy in the trap.} From the engines' 5\,s metrics, the KV held by admitted requests (SGLang's used, non-evictable
tokens, which also count locked prefixes; vLLM's cache usage) averaged over the window was 63--79\% of the pool in the six registered
SGLang 300k BURST runs, 55--83\% across all 34 stuck SGLang FCFS runs with 300k and 600k pools, 85--94\% with the native pool and 92--94\%
on vLLM, with follow-up hit fraction zero throughout. (The first computation's ``63--79\% in registered traps'' holds only for the
SGLang 300k runs.)

\paragraph{SLO goodput (D2).} Lu et al.\ define goodput as sessions completed per hour with slowdown at most 4 and at least 95\% of
decode steps meeting an inter-token latency target~\citep{lu2026cachetrap}. We have no per-token latencies, so we omit the second
condition from both arms. That makes each goodput an upper bound but does not bound the loss in either direction, except where no
BURST session meets the slowdown condition; there (26 of the 30 runs in the first hour, all 30 in the evaluation window) the loss is
100\% under either definition, provided the NOBURST goodput stays positive. Slowdown is the session's elapsed time from its planned
arrival, over an unloaded reference built from the NOBURST runs (first computation: per-turn minimum latencies; recomputation: a
least-squares fit of latency on new-suffix and decode tokens; both reported). Against the same-seed NOBURST run, base sessions that
arrived in the hour after the burst lost 99.2--100\% of goodput in all 30 registered BURST runs (no session met the target in 26), and
those arriving in the evaluation window lost 100\% in all 30, for reference latencies up to twice the unloaded ones. In the evaluation
window the CAPPED arms were within 0--1.7\% and the rescue runs within $-0.3$ to $+0.9$\% of NOBURST. In the first hour after the
burst the cap cost 4--57\% (CAPPED) and the rescue 45--97\%; these two ranges depend on the reference, and a reference twice as large
lowers them by a quarter to a third. Lu et al.\ report about 88\% from one run of a different system, which is not comparable in
detail.
